\chapter{Internal Bleeding}

\section*{Definition}
Internal bleeding is bleeding that occurs inside the body, within tissues, organs, or body cavities. It is not visible externally and can be life-threatening depending on the rate and location of blood loss.

\section*{What Happens in Your Body}
Internal bleeding results from injury to blood vessels within the body. Trauma can lacerate or rupture vessels. Non-traumatic causes include erosion of vessels by ulcers or tumors, rupture of aneurysms, bleeding from varices, or coagulopathy. Blood loss can lead to hypovolemic shock, organ compression (tamponade), and impaired organ function.

\section*{Causes}
\begin{itemize}
\item Trauma (blunt or penetrating): motor vehicle accidents, falls, assault
\item Ruptured abdominal aortic aneurysm
\item Hemorrhagic stroke (intracerebral or subarachnoid hemorrhage)
\item Peptic ulcer disease (gastrointestinal bleeding)
\item Esophageal varices (liver cirrhosis)
\item Ectopic pregnancy rupture (gynecologic emergency)
\item Ovarian cyst rupture
\item Pelvic fractures
\item Anticoagulant therapy
\item Coagulation disorders (hemophilia, liver disease)
\item Ruptured spleen (trauma, mononucleosis)
\end{itemize}

\section*{When to See a Doctor}
\begin{itemize}
\item Unexplained hypotension or shock
\item Tachycardia (rapid heart rate)
\item Pale, cool, clammy skin
\item Abdominal pain, distension, or rigidity
\item Vomiting blood (hematemesis)
\item Black, tarry stools (melena)
\item Blood in urine (hematuria)
\item Neurological deficits (if intracranial)
\item Shortness of breath (if hemothorax)
\item Confusion or altered mental status
\end{itemize}

\section*{Related Symptoms}
\begin{itemize}
\item Shock (hypovolemic)
\item Gastrointestinal bleeding
\item Hemorrhage
\item Anemia
\item Hematuria (blood in urine)
\item Hemoptysis (coughing blood)
\end{itemize}

\section*{Self-Care / Home Management}
\begin{itemize}
\item Internal bleeding is a medical emergency - CALL EMERGENCY SERVICES
\item Do NOT eat or drink anything
\item Lie down and keep calm
\item Apply direct pressure to any external bleeding sites
\item Do not move a person with suspected spinal injury
\end{itemize}

\section*{How Your Doctor Will Evaluate You}
\begin{itemize}
\item Clinical assessment of vital signs and hemodynamic status
\item Complete blood count (hemoglobin and hematocrit)
\item Coagulation profile (PT, PTT, INR)
\item CT scan of the affected area (head, chest, abdomen)
\item Focused Assessment with Sonography for Trauma (FAST) ultrasound
\item Angiography to identify bleeding source
\item Diagnostic peritoneal lavage or thoracentesis
\end{itemize}

\section*{Treatment Options}
\begin{itemize}
\item Resuscitation: IV fluids, blood transfusion
\item Reverse anticoagulation if applicable
\item Surgical intervention to control bleeding source
\item Endoscopic therapy for GI bleeding
\item Angiographic embolization
\item Treat the underlying cause
\item ICU monitoring for hemorrhagic shock
\end{itemize}

\section*{References}
\begin{thebibliography}{99}
\bibitem{ref1} Kauvar DS, Lefering R, Wade CE. "Impact of hemorrhage on trauma outcome: an overview of epidemiology, clinical presentations, and therapeutic considerations." \textit{Journal of Trauma}, 2006;60(6):S3--S11.
\bibitem{ref2} Strate LL, Gralnek IM. "ACG clinical guideline: management of patients with acute lower gastrointestinal bleeding." \textit{American Journal of Gastroenterology}, 2016;111(4):459--474.
\bibitem{ref3} Barkun AN, Bardou M, Kuipers EJ, et al. "International consensus recommendations on the management of patients with nonvariceal upper gastrointestinal bleeding." \textit{Annals of Internal Medicine}, 2010;152(2):101--113.
\bibitem{ref4} Spahn DR, Bouillon B, Cerny V, et al. "The European guideline on management of major bleeding and coagulopathy following trauma." \textit{Critical Care}, 2019;23(1):98.
\bibitem{ref5} Guirguis-Blake JM, Evans CV, Rushkin M, et al. "Screening for abdominal aortic aneurysm." \textit{JAMA}, 2019;322(22):2212--2218.
\bibitem{ref6} Laine L, Jensen DM. "Management of patients with ulcer bleeding." \textit{American Journal of Gastroenterology}, 2012;107(3):345--360.

\end{thebibliography}

\clearpage
